% BHU PYQ -- Sociology: Understanding of Indian Society (2024-25 NEP)
\documentclass[11pt,a4paper]{article}
\usepackage[a4paper,top=2.5cm,bottom=2.5cm,left=2.8cm,right=2.8cm,headheight=14pt]{geometry}
\usepackage{amsmath,amssymb}


\usepackage{fontspec}
\setmainfont{Kohinoor Devanagari}
\usepackage{microtype,fancyhdr,enumitem,hyperref,lastpage,parskip}

\hypersetup{colorlinks=true,urlcolor=black,linkcolor=black,pdftitle={SOCMJ21: Understanding of Indian Society -- BHU NEP 2024-25}}

\pagestyle{fancy}\fancyhf{}
\renewcommand{\headrulewidth}{0.4pt}\renewcommand{\footrulewidth}{0.4pt}
\fancyhead[L]{\small \textit{Banaras Hindu University}}
\fancyhead[R]{\small \textit{SOCMJ21: Indian Society (2024-25)}}
\fancyfoot[C]{\small Page \thepage\ of \pageref{LastPage}}

\newlist{parts}{enumerate}{2}
\setlist[parts,1]{label=(\alph*),leftmargin=*,itemsep=4pt,topsep=3pt}
\setlist[parts,2]{label=(\roman*),leftmargin=*,itemsep=2pt,topsep=2pt}

\newcommand{\pts}[1]{\hfill \textit{[#1]}}

\begin{document}

\begin{center}
  {\large\bfseries BANARAS HINDU UNIVERSITY}\\[2pt]
  {\normalsize B. A. (Hons.) Semester II Examination, 2024-25 (NEP)}\\[6pt]
  \rule{0.8\linewidth}{0.5pt}\\[6pt]
  {\large\bfseries SOCIOLOGY}\\[4pt]
  {\large\bfseries Paper Code: SOCMJ-21 : Understanding of Indian Society}\\[4pt]
  {\normalsize Time: 2:30 Hours \hfill Full Marks: 70}\\[2pt]
  {\small REF NO. D-II/077}
\end{center}
\rule{\linewidth}{0.4pt}
\smallskip

\noindent \textit{Instructions: This question paper is divided into three Sections A, B and C. Attempt all sections according to directions.}

\smallskip
\noindent \textit{यह प्रश्न-पत्र तीन खण्डों अ, ब तथा स में विभाजित है। निर्देशों के अनुसार सभी खण्डों को हल कीजिए।}

\bigskip

\section*{SECTION -- A / खण्ड -- अ}
\noindent \textit{Note: Very Short Answer Type Questions. Each question carries 2 marks. Answer within 50 words.}\pts{5 $\times$ 2 = 10}

\noindent \textit{अति लघु उत्तरीय प्रश्न। प्रत्येक प्रश्न 2 अंकों का है। उत्तर 50 शब्दों से अधिक नहीं होना चाहिए।}

\begin{enumerate}[label=\textbf{\arabic*.},leftmargin=*,itemsep=8pt]
  \item 
  \begin{parts}
    \item Who said `Caste is a closed class'?\pts{2}
    
    `जाति एक बंद वर्ग है' किसने कहा है? (D. N. Majumdar)

    \item Who is the author of the book `Social Change in Modern India'?\pts{2}
    
    `सोशल चेंज इन मॉडर्न इण्डिया' पुस्तक के लेखक कौन हैं? (M. N. Srinivas)

    \item What do you mean by Kinship?\pts{2}
    
    नातेदारी से आप क्या समझते हैं?

    \item Who is the author of the book `Social Background of Indian Nationalism'?\pts{2}
    
    `सोशल बैकग्राउंड ऑफ इण्डियन नेशनलिज्म' पुस्तक के लेखक कौन हैं? (A. R. Desai)

    \item What is hypergamy and hypogamy marriage?\pts{2}
    
    अनुलोम तथा प्रतिलोम विवाह क्या हैं?
  \end{parts}
\end{enumerate}

\bigskip

\section*{SECTION -- B / खण्ड -- ब}
\noindent \textit{Note: Short Answer Type Questions. Each question carries 10 marks. Answer within 250 words.}\pts{3 $\times$ 10 = 30}

\noindent \textit{लघु उत्तरीय प्रश्न। प्रत्येक प्रश्न 10 अंकों का है। उत्तर 250 शब्दों से अधिक नहीं होना चाहिए।}

\begin{enumerate}[label=\textbf{\arabic*.},leftmargin=*,start=2,itemsep=12pt]

  \item Write a short note on the Nationalist discourse of understanding Indian Society.\pts{10}
  
  भारतीय समाज को समझने के राष्ट्रवादी परिप्रेक्ष्य पर एक संक्षिप्त टिप्पणी लिखिए।

  \begin{center}\textbf{OR / अथवा}\end{center}

  What is Caste? Discuss the basic characteristics of the caste system in India.\pts{10}
  
  जाति क्या है? भारत में जाति व्यवस्था की मौलिक विशेषताओं की चर्चा कीजिए।

  \item What do you understand by Tribal Society? Explain its salient features.\pts{10}
  
  जनजातीय समाज से आप क्या समझते हैं? इसकी प्रमुख विशेषताओं की व्याख्या कीजिए।

  \begin{center}\textbf{OR / अथवा}\end{center}

  Clarify the concept of Class. Differentiate between Class and Caste.\pts{10}
  
  वर्ग की अवधारणा को स्पष्ट कीजिए। वर्ग एवं जाति में अन्तर स्पष्ट कीजिए।

  \item Discuss different types of kinship usages in India.\pts{10}
  
  भारत में विभिन्न प्रकार की नातेदारी की रीतियों (व्यवहारों) की व्याख्या कीजिए।

  \begin{center}\textbf{OR / अथवा}\end{center}

  Write a note on the emergence of class formation in India.\pts{10}
  
  भारत में वर्ग निर्माण के उदय पर एक टिप्पणी लिखिए।

\end{enumerate}

\bigskip

\section*{SECTION -- C / खण्ड -- स}
\noindent \textit{Note: Long Answer Type Questions. Each question carries 15 marks. Answer within 500 words.}\pts{2 $\times$ 15 = 30}

\noindent \textit{दीर्घ उत्तरीय प्रश्न। प्रत्येक प्रश्न 15 अंकों का है। उत्तर 500 शब्दों से अधिक नहीं होना चाहिए।}

\begin{enumerate}[label=\textbf{\arabic*.},leftmargin=*,start=5,itemsep=14pt]

  \item How does colonial discourse view Indian Society? Explain in detail.\pts{15}
  
  औपनिवेशिक परिप्रेक्ष्य भारतीय समाज को किस रूप में देखता है? व्याख्या कीजिए।

  \begin{center}\textbf{OR / अथवा}\end{center}

  Write an essay on the changing patterns of family in India.\pts{15}
  
  भारत में परिवार के बदलते प्रतिमानों पर एक निबन्ध लिखिए।

  \item Discuss the different types of Marriage in Hindu Tradition.\pts{15}
  
  हिन्दू परम्परा में विवाह के विभिन्न प्रकारों की चर्चा कीजिए।

  \begin{center}\textbf{OR / अथवा}\end{center}

  Discuss the developmental challenges faced by Tribal Society in India.\pts{15}
  
  भारत में जनजातीय समाज द्वारा सामना की जा रही विकास सम्बन्धी चुनौतियों की चर्चा कीजिए।

\end{enumerate}

\end{document}
